\documentclass[11pt]{article}
\usepackage[a4paper,margin=22mm]{geometry}
\usepackage{fontspec}
\setmainfont{DejaVu Serif}
\setsansfont{DejaVu Sans}
\usepackage{microtype}
\usepackage{xcolor}
\usepackage{hyperref}
\usepackage{enumitem}
\usepackage{parskip}
\definecolor{Accent}{HTML}{971D32}
\definecolor{Muted}{HTML}{666666}
\hypersetup{colorlinks=true,urlcolor=Accent,linkcolor=Accent}
\setlist{leftmargin=1.6em,itemsep=0.3em,topsep=0.35em}
\setcounter{tocdepth}{2}
\renewcommand{\contentsname}{Contents}
\newcommand{\sectionrule}{\par\vspace{-0.5em}\noindent\textcolor{Accent}{\rule{\linewidth}{0.6pt}}\par}
\begin{document}

\begin{center}
{\sffamily\bfseries\color{Accent} TOEFL WORD OF THE DAY\par}
\vspace{8mm}
{\fontsize{42}{48}\selectfont\bfseries propel\par}
\vspace{2mm}
{\Large\sffamily\color{Accent} /prəˈpɛl/\par}
\vspace{2mm}
{\small\sffamily Local audio phrase: propel\par}
{\small\sffamily Audio file: audios/propel.mp3\par}
\vspace{2mm}
{\large\sffamily\color{Muted} transitive verb\par}
\vspace{5mm}
{\sffamily push or drive something forward \textbullet\ cause rapid progress or advancement\par}
\vspace{6mm}
{\large Propel means to drive something forward physically or to cause rapid progress toward a new position or situation.\par}
\vspace{6mm}
\begin{minipage}{0.84\textwidth}
\itshape “Her discovery propelled her to international recognition.”
\end{minipage}\par
\vspace{10mm}
\textbf{Date:} August 14th 2026\quad\textbullet\quad \textbf{Level:} B1--mid-B2 TOEFL
\vfill
Created and compiled by \textbf{Amir Kooshky}\par
\vspace{2mm}
\href{https://t.me/KooshkyTOEFL}{Telegram Channel}\quad\textbullet\quad
\href{https://www.instagram.com/kooshkytoefl}{Instagram}\quad\textbullet\quad
\href{https://t.me/KooshkyTOEFL\_pv}{Personal Telegram}
\end{center}
\newpage
\tableofcontents
\newpage

\section{Meanings and usage}
\sectionrule
\subsection{Meaning 1: Push or drive forward}
\textbf{Part of speech:} transitive verb

\textbf{IPA:} /prəˈpɛl/

\textbf{Pronunciation phrase:} propel

\textbf{Local audio file:} audios/propel.mp3

\textbf{Register:} neutral to formal; common in science, engineering, biology, transportation, and descriptions of movement

\textbf{Definition:} to push, drive, or cause a person or object to move forward or in a particular direction

\textbf{In simpler words:} push or drive forward

\textbf{Typical pattern:} propel + object + forward, through, toward, or into + place

\subsubsection{Natural examples}
\begin{enumerate}
  \item Powerful engines propel the aircraft through the air.
  \item The squid propels itself through the water by forcing water out of its body.
  \item A strong current propelled the boat toward the shore.
  \item Gas released from the rocket engine propels the spacecraft forward.
  \item The swimmer used her legs to propel herself across the pool.
  \item Wind can propel small seeds over considerable distances.
  \item The explosion propelled pieces of rock high into the air.
  \item Some marine animals propel themselves by rapidly moving their tails.
\end{enumerate}

\subsubsection{Nuance and implication}
\textbf{Central nuance:} Propel emphasizes the force or mechanism that produces movement, especially forward movement.

\textbf{Usual implication:} The subject supplies energy or force that causes another object, or itself, to move.

\textbf{Compared with carry:} Carry means support and transport something from one place to another. Propel emphasizes the force that makes something move.

\subsubsection{Grammar notes}
\textbf{How to use it:} This verb normally needs the person or thing being moved directly after it.

\textbf{What can follow it:} After the object, you can add a direction or destination with words such as forward, through, toward, across, or into.

\textbf{Common preposition:} Say propel something through the water, toward the shore, or into the air.

\textbf{Position in a sentence:} When the subject moves itself, use a reflexive form such as propel itself or propel herself.

\subsubsection{Useful patterns}
\textbf{Pattern:} propel + object + forward

\textbf{Use:} Use this when a force or mechanism causes something to move ahead.

\textbf{Example:} The motor propels the vehicle forward.

\textbf{Pattern:} propel + object + through + substance or space

\textbf{Use:} Use this to describe movement through water, air, or another environment.

\textbf{Example:} The fins help propel the animal through the water.

\textbf{Pattern:} propel + object + toward + destination

\textbf{Use:} Use this when movement is directed toward a particular place.

\textbf{Example:} The current propelled the raft toward the waterfall.

\textbf{Pattern:} propel oneself + direction

\textbf{Use:} Use this when a person or animal creates the force for its own movement.

\textbf{Example:} The diver propelled herself upward with a powerful kick.

\textbf{Pattern:} be propelled by + force or mechanism

\textbf{Use:} Use the passive form when emphasizing what provides the force.

\textbf{Example:} The vessel is propelled by two electric motors.

\subsubsection{Common wrong usage}
\paragraph{Correction 1}
\textbf{Wrong:} The engine propels forward the boat.

\textbf{Correct:} The engine propels the boat forward.

\textbf{Why:} Place the object directly after propel, and put forward after the object.

\paragraph{Correction 2}
\textbf{Wrong:} The fish propels through the water with its tail.

\textbf{Correct:} The fish propels itself through the water with its tail.

\textbf{Why:} Propel is normally transitive, so when the subject causes its own movement, use a reflexive pronoun such as itself.

\paragraph{Correction 3}
\textbf{Wrong:} The rocket was propelled from powerful engines.

\textbf{Correct:} The rocket was propelled by powerful engines.

\textbf{Why:} Use by to identify the force or mechanism that causes the movement.

\subsection{Meaning 2: Cause rapid progress or advancement}
\textbf{Part of speech:} transitive verb

\textbf{IPA:} /prəˈpɛl/

\textbf{Pronunciation phrase:} propel

\textbf{Local audio file:} audios/propel.mp3

\textbf{Register:} neutral to formal; common in academic, business, political, historical, and journalistic writing

\textbf{Definition:} to cause a person, organization, idea, or process to advance quickly or reach a new position, level, or situation

\textbf{In simpler words:} make someone or something advance quickly

\textbf{Typical pattern:} propel + person or thing + to, into, or toward + new position or situation

\subsubsection{Natural examples}
\begin{enumerate}
  \item Her discovery propelled her to international recognition.
  \item Rapid technological change has propelled the industry into a period of intense competition.
  \item Increased demand helped propel the company's growth during its first five years.
  \item The successful campaign propelled the candidate into national politics.
  \item A series of influential studies propelled the theory to the center of scientific debate.
  \item Strong export growth helped propel the economy out of recession.
  \item The novel's popularity propelled its author to fame.
  \item New investment could propel the region toward faster economic development.
\end{enumerate}

\subsubsection{Nuance and implication}
\textbf{Central nuance:} This figurative sense presents progress or change as if a force were pushing someone or something forward.

\textbf{Usual implication:} The change is often significant and relatively rapid, moving the person, organization, or idea into a noticeably different position.

\textbf{Compared with promote:} Promote can mean actively support or help something develop. Propel emphasizes the strong effect that causes rapid advancement or movement into a new position.

\subsubsection{Grammar notes}
\textbf{How to use it:} Place the person, organization, idea, or process directly after propel.

\textbf{What can follow it:} Use to for a new status or level, into for a new situation or field, and toward for movement in the direction of a goal.

\textbf{Common preposition:} Common patterns include propel someone to fame, propel someone into leadership, and propel a country toward recovery.

\textbf{Position in a sentence:} The passive form is also common, as in She was propelled to fame by the success of her first novel.

\subsubsection{Useful patterns}
\textbf{Pattern:} propel + person + to + fame, success, or prominence

\textbf{Use:} Use this when an event causes someone's status to rise rapidly.

\textbf{Example:} The award propelled the young scientist to international prominence.

\textbf{Pattern:} propel + person or organization + into + position or situation

\textbf{Use:} Use this when something causes entry into a new role, field, or condition.

\textbf{Example:} The election propelled the movement into national politics.

\textbf{Pattern:} propel + thing + toward + goal or change

\textbf{Use:} Use this for rapid movement in the direction of a desired or expected result.

\textbf{Example:} New investment could propel the country toward economic recovery.

\textbf{Pattern:} propel growth or development

\textbf{Use:} Use this when a factor strongly accelerates an economic, social, or organizational process.

\textbf{Example:} Rising demand propelled the growth of the renewable energy industry.

\textbf{Pattern:} be propelled by + event or factor

\textbf{Use:} Use this when emphasizing what caused the rapid advancement.

\textbf{Example:} The company's expansion was propelled by strong overseas demand.

\subsubsection{Common wrong usage}
\paragraph{Correction 1}
\textbf{Wrong:} The discovery propelled her for international fame.

\textbf{Correct:} The discovery propelled her to international fame.

\textbf{Why:} Use propel someone to a new level, status, or achievement.

\paragraph{Correction 2}
\textbf{Wrong:} The new investment propelled to the company's growth.

\textbf{Correct:} The new investment propelled the company's growth.

\textbf{Why:} Propel normally takes the thing being advanced directly as its object.

\paragraph{Correction 3}
\textbf{Wrong:} The election propelled him in national politics.

\textbf{Correct:} The election propelled him into national politics.

\textbf{Why:} Use into when someone enters a new field, role, or situation.

\paragraph{Correction 4}
\textbf{Wrong:} The policy propelled the economy to recover.

\textbf{Correct:} The policy helped propel the economy toward recovery.

\textbf{Why:} For movement in the direction of a process or goal, propel something toward a noun such as recovery is more natural.

\section{Word family}
\sectionrule
\subsection{propulsion}
\textbf{Part of speech:} uncountable noun

\textbf{IPA:} /prəˈpʌlʃən/

\textbf{Definition:} the force or system that causes a vehicle or object to move forward

\textbf{Relationship to the headword:} This noun refers to the force, process, or technology used to propel something.

\textbf{Grammar pattern:} propulsion system, method of propulsion, or propulsion technology

\textbf{Usage note:} It is especially common in engineering, transportation, aerospace, and physics.

\subsubsection{Examples}
\begin{enumerate}
  \item The spacecraft uses an electric propulsion system.
  \item Engineers are developing more efficient methods of marine propulsion.
  \item Jet propulsion made much faster aircraft possible.
\end{enumerate}

\subsection{propellant}
\textbf{Part of speech:} countable or uncountable noun

\textbf{IPA:} /prəˈpɛlənt/

\textbf{Definition:} a substance used to produce the force that moves a rocket, projectile, or other object

\textbf{Relationship to the headword:} A propellant is a substance that produces the force needed to propel something.

\textbf{Grammar pattern:} rocket propellant, solid propellant, or liquid propellant

\textbf{Usage note:} This is mainly a scientific and technical term.

\subsubsection{Examples}
\begin{enumerate}
  \item The rocket carries both fuel and an oxidizer as propellants.
  \item Scientists tested a new propellant for use in small spacecraft.
\end{enumerate}

\subsection{propeller}
\textbf{Part of speech:} countable noun

\textbf{IPA:} /prəˈpɛlər/

\textbf{Definition:} a device with rotating blades that pushes against air or water to move a vehicle forward

\textbf{Relationship to the headword:} A propeller is a device that propels a boat, aircraft, or other vehicle.

\textbf{Grammar pattern:} boat, aircraft, or ship + propeller

\subsubsection{Examples}
\begin{enumerate}
  \item The boat's propeller was damaged by underwater debris.
  \item Each aircraft engine turns a large propeller.
\end{enumerate}

\section{Common collocations}
\sectionrule
\subsection{propel forward}
\textbf{Applies to:} Push or drive forward

\textbf{Meaning:} cause something to move ahead

\textbf{Pattern:} propel + object + forward

\textbf{Example:} The rear wheels propel the vehicle forward.

\subsection{propel through the water}
\textbf{Applies to:} Push or drive forward

\textbf{Meaning:} cause an animal or object to move through water

\textbf{Pattern:} propel + object + through the water

\textbf{Example:} Powerful tail movements propel the whale through the water.

\subsection{propel itself}
\textbf{Applies to:} Push or drive forward

\textbf{Meaning:} cause one's own body or structure to move

\textbf{Pattern:} propel itself, himself, herself, or themselves + direction

\textbf{Example:} The squid propels itself by rapidly expelling water.

\subsection{be propelled by}
\textbf{Applies to:} Push or drive forward

\textbf{Meaning:} be moved by a particular force or mechanism

\textbf{Pattern:} be propelled by + force or mechanism

\textbf{Example:} The submarine is propelled by an electric motor.

\subsection{propel to fame}
\textbf{Applies to:} Cause rapid progress or advancement

\textbf{Meaning:} cause someone to become famous very quickly

\textbf{Pattern:} propel + person + to fame

\textbf{Example:} The unexpected success of the film propelled the actor to fame.

\subsection{propel to prominence}
\textbf{Applies to:} Cause rapid progress or advancement

\textbf{Meaning:} cause someone to become important or well known

\textbf{Pattern:} propel + person + to prominence

\textbf{Example:} The discovery propelled the researcher to international prominence.

\subsection{propel into}
\textbf{Applies to:} Cause rapid progress or advancement

\textbf{Meaning:} cause someone or something to enter a new position or situation

\textbf{Pattern:} propel + person or thing + into + situation

\textbf{Example:} The crisis propelled the issue into the center of public debate.

\subsection{propel growth}
\textbf{Applies to:} Cause rapid progress or advancement

\textbf{Meaning:} cause growth to occur more rapidly

\textbf{Example:} Low interest rates helped propel economic growth.

\subsection{propel development}
\textbf{Applies to:} Cause rapid progress or advancement

\textbf{Meaning:} cause development to advance more rapidly

\textbf{Example:} Investment in transportation infrastructure propelled the development of the region.

\subsection{propel toward}
\textbf{Applies to:} Cause rapid progress or advancement

\textbf{Meaning:} cause movement in the direction of a new condition or goal

\textbf{Pattern:} propel + person or thing + toward + goal

\textbf{Example:} The reforms could propel the country toward greater economic stability.

\subsection{propulsion system}
\textbf{Applies to:} Push or drive forward

\textbf{Meaning:} a system that provides the force necessary for movement

\textbf{Example:} The spacecraft uses a highly efficient propulsion system.

\textbf{Usage note:} Propulsion is the noun form.

\section{Synonyms and antonyms}
\sectionrule
\subsection{Close synonyms}
\subsubsection{drive}
\textbf{Part of speech:} transitive verb

\textbf{Applies to:} Push or drive forward

\textbf{Shared meaning:} Both can describe a force causing movement.

\textbf{Difference:} Drive is broader and can mean make something move, operate a vehicle, or strongly motivate change. Propel specifically emphasizes producing forward movement.

\textbf{Example:} The strong wind drove the boat toward the coast.

\subsubsection{push}
\textbf{Part of speech:} transitive verb

\textbf{Applies to:} Push or drive forward

\textbf{Shared meaning:} Both can mean apply force that causes movement.

\textbf{Difference:} Push is the ordinary general word for applying force to move something away from you. Propel is more formal and often describes continuous or powerful forward movement.

\textbf{Example:} The workers pushed the cart across the warehouse.

\subsubsection{accelerate}
\textbf{Part of speech:} transitive verb

\textbf{Applies to:} Cause rapid progress or advancement

\textbf{Shared meaning:} Both can describe making progress occur more rapidly.

\textbf{Difference:} Accelerate means make a process happen faster. Propel suggests both speed and movement toward a new stage, position, or outcome.

\textbf{Example:} New technology accelerated the development of renewable energy systems.

\subsubsection{advance}
\textbf{Part of speech:} transitive verb

\textbf{Applies to:} Cause rapid progress or advancement

\textbf{Shared meaning:} Both can refer to causing progress.

\textbf{Difference:} Advance means help something develop or move forward. Propel is stronger and often suggests a powerful factor causing rapid or dramatic progress.

\textbf{Example:} The discovery advanced scientists' understanding of the disease.

\section{Words learners may confuse}
\sectionrule
\subsection{compel}
\textbf{Part of speech:} transitive verb

\textbf{Why learners confuse them:} The words have similar spellings and both can suggest a strong force, but the type of force is different.

\textbf{Key difference:} Propel means push or cause movement or progress. Compel means force someone to do something.

\textbf{propel:} Strong demand propelled the company into new markets.

\textbf{compel:} The law compels employers to provide basic safety equipment.

\subsection{propagate}
\textbf{Part of speech:} verb

\textbf{Why learners confuse them:} Both formal verbs may appear in scientific descriptions of movement, but they describe different processes.

\textbf{Key difference:} Propel means cause physical movement or rapid advancement. Propagate means spread an idea, belief, signal, organism, or other thing.

\textbf{propel:} The engine propels the aircraft forward.

\textbf{propagate:} Radio waves propagate through the atmosphere.

\section{Etymology}
\sectionrule
\textbf{Origin:} Propel comes from Latin propellere, meaning to drive or push forward.

\textbf{Historical form:} Latin propellere

\textbf{Early meaning:} to drive forward or push away

\textbf{Memory link:} The idea of pushing forward remains central in both the physical and figurative meanings of propel.

\section{Practice exercises}
\sectionrule
\subsection{Word family choice}
Choose the best member of the word family for each blank.

\begin{enumerate}
  \item The spacecraft's new \_\_\_\_\_ system uses electric power rather than conventional chemical rockets.
  \item The rocket requires a specially designed \_\_\_\_\_ to generate thrust.
  \item The boat's \_\_\_\_\_ converts engine power into forward motion.
  \item The new evidence could \_\_\_\_\_ the theory to greater prominence.
\end{enumerate}

\subsection{Correct or incorrect?}
Decide whether each sentence is correct. Correct the inaccurate ones.

\begin{enumerate}
  \item The engine propels the boat forward.
  \item The fish propels through the water by moving its tail.
  \item The successful novel propelled the author to international fame.
  \item The election propelled her in national politics.
  \item Strong consumer demand propelled the company's growth.
\end{enumerate}

\subsection{Collocation practice}
Complete each sentence with a natural collocation.

\begin{enumerate}
  \item The vessel is propelled \_\_\_\_\_ two powerful electric motors.
  \item The award propelled the young researcher \_\_\_\_\_ international prominence.
  \item The squid propels \_\_\_\_\_ through the water by expelling water from its body.
  \item Rapid investment helped \_\_\_\_\_ the region's economic growth.
  \item The unexpected victory propelled the movement \_\_\_\_\_ national politics.
\end{enumerate}

\subsection{Complete answer key}
\subsubsection{Word family choice}
\begin{enumerate}
  \item \textbf{propulsion.} The spacecraft's new propulsion system uses electric power rather than conventional chemical rockets. \emph{Use propulsion for the system or force that moves a vehicle.}
  \item \textbf{propellant.} The rocket requires a specially designed propellant to generate thrust. \emph{A propellant is a substance used to produce the force needed for movement.}
  \item \textbf{propeller.} The boat's propeller converts engine power into forward motion. \emph{A propeller is the rotating device that helps propel a boat or aircraft.}
  \item \textbf{propel.} The new evidence could propel the theory to greater prominence. \emph{Use the base verb propel after could.}
\end{enumerate}

\subsubsection{Correct or incorrect?}
\begin{enumerate}
  \item \textbf{Correct} \emph{Propel the boat forward is a standard physical use of the verb.}
  \item \textbf{Incorrect}. The fish propels itself through the water by moving its tail. \emph{Propel normally needs an object, so use itself when the animal causes its own movement.}
  \item \textbf{Correct} \emph{Propel someone to fame is a common figurative pattern.}
  \item \textbf{Incorrect}. The election propelled her into national politics. \emph{Use into when someone is moved into a new field, role, or situation.}
  \item \textbf{Correct} \emph{Propel growth naturally means cause growth to advance rapidly.}
\end{enumerate}

\subsubsection{Collocation practice}
\begin{enumerate}
  \item \textbf{by.} The vessel is propelled by two powerful electric motors. \emph{Use by to identify the force or mechanism that causes movement.}
  \item \textbf{to.} The award propelled the young researcher to international prominence. \emph{Propel someone to prominence means cause them to become important or well known quickly.}
  \item \textbf{itself.} The squid propels itself through the water by expelling water from its body. \emph{Use a reflexive pronoun when an animal or person causes its own movement.}
  \item \textbf{propel.} Rapid investment helped propel the region's economic growth. \emph{Propel growth means cause growth to advance quickly.}
  \item \textbf{into.} The unexpected victory propelled the movement into national politics. \emph{Propel something into a new situation means cause it to enter that situation rapidly.}
\end{enumerate}

\vfill
\begin{center}
\small\color{Muted} Created and compiled by \textbf{Amir Kooshky}\\
\href{https://t.me/KooshkyTOEFL}{Telegram Channel}\quad\textbullet\quad
\href{https://www.instagram.com/kooshkytoefl}{Instagram}\quad\textbullet\quad
\href{https://t.me/KooshkyTOEFL\_pv}{Personal Telegram}
\end{center}
\end{document}
